\section{Results on the Benchmark}
\label{sec:results}
This section reports the 68 benchmark cases (\NRunsBenchmark{} runs of 6,030 evaluations, random starts, 15$^\circ$ rose; per case: Tables~\ref{S-tab:res-1-500}--\ref{S-tab:res-2-1000}).

\subsection{Effect of the Baseline Setting}
\label{sec:psosetting}
\label{sec:baseline}
\IfFileExists{analysis/mpce_tab_baseline.tex}{%% generated by mpce_results.py -- do not edit by hand
\begin{table}[!t]
\centering
\caption{Effect of the PSO Setting on the Comparison of the Previous Study's Methods (68 Cases, 6,030 Evaluations): Average Rank (1 = Best) with the Old ($w=0.7$, $c_1=c_2=2$) and the Constriction Setting, and Run-Level W/T/L of the Salp-Swarm Hybrids against PSO}
\label{tab:baseline}
\footnotesize\setlength{\tabcolsep}{4pt}
\begin{tabular}{lcc}
\toprule
& \multicolumn{2}{c}{PSO setting} \\
\cmidrule(lr){2-3}
Method & Old & Constriction \\
\midrule
LX-SSA-VNS & 2.65 & 3.41 \\
SSA-VNS & \textbf{1.99} & 2.77 \\
LX-SSA & 5.65 & 6.06 \\
SSA & 4.81 & 5.34 \\
PSO & 5.04 & \textbf{1.53} \\
DE & 6.99 & 7.03 \\
VNS & 3.59 & 4.15 \\
MS-SLSQP & 5.28 & 5.71 \\
\midrule
PSO position & 5 & 1 \\
LX-SSA-VNS vs.\ PSO (W/T/L) & 35/33/0 & 0/21/47 \\
SSA-VNS vs.\ PSO (W/T/L) & 41/27/0 & 0/22/46 \\
\bottomrule
\multicolumn{3}{p{0.95\columnwidth}}{Ranks: feasibility-aware rule (fewer than 15 of 30 feasible runs in a case = ranked last); bold: best average rank. W/T/L: cases in which the hybrid is significantly better / not different / worse than PSO (two-sided Wilcoxon signed-rank test, 30 seed-paired runs, Holm-adjusted over the seven comparisons of the hybrid in each case, $\alpha=0.05$).}
\end{tabular}
\end{table}
}{}
The old setting lies beyond the order-2 (variance) stability boundary~\eqref{eq:poli}. Run on the same cases, budget and seeds as the constriction setting~\eqref{eq:constriction}, it returns a feasible layout in only \NOldPSOFeas\% of the runs (fewer than 15 feasible runs in \NOldPSOBelowHalf{} cases), is significantly worse in \NOldPSOWorse{} of the 68 cases and never better, and has a mean wake loss \NOldPSODL{}~pp higher over the \NOldPSOCases{} cases in which both settings have at least 15 feasible runs. In instrumented reruns (clipping without velocity reset or clamping; Table~\ref{S-tab:D-psodyn}), its swarm contracts much less (final spread \NDSpreadRatio{} times, velocity \NDVelRatio{} times that of the constriction setting), \NDClipOld\% of the coordinates are clipped (constriction: \NDClipNew\%), and far fewer particle evaluations in the second half of a run are feasible (\NDFeasPartOld\% vs.\ \NDFeasPartNew\%). A sweep of $c_1=c_2=c$ at $w=0.7$ (Fig.~\ref{S-fig:S-csweep}) shows that the boundary $c^\ast=\NSBoundC$ marks approximately where this PSO starts to fail: at least \NSFeasBelow\% of the runs are feasible for $c\le\NSCBelow$, then feasibility declines progressively (\NSFeasAbove\% at $c=\NSCAbove$, \NSFeasCTwoZero\% at 2.0); velocity clamping raises it to \NSFeasVMax\%, still worse than constriction in all \NSLossCasesVMaxConstr{} cases.
% CHECK-FINAL [C12]: pso_setting: constriction_vs_old L == 0, W >= 34 (editor cut C5 applied)
% CHECK-DIAG [D03] spread/|V| ratios ("contracts much less"); [D04] clipped share; [D05] feasible particle evaluations (iterations 101-200); 
% CHECK-CSWEEP [S02] c <= 1.7 >= 98 % feasible; [S03] c = 1.8 > 1.9 > 2.0; [S04] decline starts in the step containing c*; [S05] progressive, no jump; [S10] vmax >= 90 %; [S11] vmax worse than constriction in all 12 cases (never "causes")

Table~\ref{tab:baseline} shows the effect on a comparison. It ranks an eight-method pool, the methods of Table~\ref{tab:friedman68} with LX-SSA-VNS instead of PSO-VNS, once with each PSO setting (all other methods use identical runs): with the old setting, \NBaseOldBest{} ranks first and PSO only in position \NBaseOldPSOPos{}, behind both salp-swarm hybrids; with the constriction setting, \NBaseNewBest{} ranks first, ahead of every salp-swarm method (standard coefficients~\cite{Clerc2002}; none of the methods was tuned on these cases). The reversal holds at 6,030 evaluations; on the six largest cases at 120,030, SSA-VNS ranks ahead of PSO (\NBudRankSSAVNSOneTwentyK{} vs.\ \NBudRankPSOOneTwentyK; Section~\ref{sec:budget}). A PSO baseline beyond the order-2 boundary thus makes methods compared with it, including our own LX-SSA, appear stronger than they are.
% CHECK-FINAL [C27]: same pool: old setting SSA-VNS (1.99) and LX-SSA-VNS (2.65) ahead of PSO (5.04, NBaseOldPSOPos); constriction PSO ahead of every salp-swarm method (NBaseNewBest = PSO)
% CHECK-FINAL [C45]: at 120,030 evaluations SSA-VNS ranks ahead of PSO on the six largest cases

\subsection{Overall Comparison}
\label{sec:overall}
%% generated by mpce_results.py -- do not edit by hand
\begin{table*}[!t]
\centering
\caption{Case-level analysis over the 68 benchmark cases. Ranks of the mean feasible objective within each case (1 = best); a method with fewer than 15 feasible runs out of 30 in a case is ranked below all other methods, by its number of feasible runs (``$<$15'': number of such cases). Friedman $\chi^2_F=174.4$ (6 d.f.), $p=5.29\times10^{-35}$; Iman--Davenport $F_F=50.0$, $p=8.06\times10^{-46}$. $p_z$: Holm-adjusted $p$ of the average-rank $z$ test against SSA-VNS (best-ranked method: SSA-VNS). Because mean-rank post hoc tests depend on the pool of compared methods \cite{Benavoli2016}, pairwise two-sided Wilcoxon signed-rank tests on the per-case mean wake losses (\%) are also given: $p_W$ (Holm-adjusted over the 6 comparisons), matched-pairs rank-biserial correlation $r_{\rm rb}$ (positive = SSA-VNS better) and mean difference $\overline{\Delta L}$ of the wake loss (percentage points, SSA-VNS minus method, over the cases where both methods have at least 15 feasible runs). ``Best'': cases in which the method alone ranks first; ``Feas.'': percentage of feasible runs.}
\label{tab:friedman68}
\scriptsize\setlength{\tabcolsep}{3pt}
\begin{tabular}{lcccccccc}
\toprule
Method & Avg.\ rank & Best & Feas.\ (\%) & $<$15 & $p_z$ & $p_W$ & $r_{\rm rb}$ & $\overline{\Delta L}$ (pp) \\
\midrule
SSA-VNS & 1.74 & 37 & 99.6 & 0 & -- & -- & -- & -- \\
VNS & 2.97 & 8 & 100.0 & 0 & $8.55\times10^{-4}$ & $5.49\times10^{-7}$ & +0.77 & $-0.135$ \\
SSA & 3.88 & 0 & 97.6 & 2 & $1.36\times10^{-8}$ & $2.27\times10^{-10}$ & +0.96 & $-0.294$ \\
PSO & 4.17 & 3 & 86.0 & 8 & $1.52\times10^{-10}$ & $1.44\times10^{-9}$ & +0.92 & $-0.337$ \\
MS-SLSQP & 4.47 & 11 & 99.9 & 0 & $6.18\times10^{-13}$ & $3.34\times10^{-5}$ & +0.58 & $-0.112$ \\
LX-SSA & 4.69 & 0 & 96.9 & 2 & $7.40\times10^{-15}$ & $5.53\times10^{-11}$ & +1.00 & $-0.447$ \\
DE & 6.08 & 0 & 71.4 & 18 & $5.39\times10^{-31}$ & $7.74\times10^{-11}$ & +1.00 & $-0.680$ \\
\bottomrule
\end{tabular}
\end{table*}

%% generated by mpce_results.py -- do not edit by hand
\begin{table}[!t]
\centering
\caption{Pairwise outcome of PSO-VNS against each method: number of cases in which PSO-VNS is significantly better / not significantly different / significantly worse (two-sided Wilcoxon signed-rank test on 30 seed-paired runs, Holm-adjusted over the 7 comparisons of each case, $\alpha=0.05$; infeasible runs rank below all feasible runs). Last row: median over the cases of the rank-biserial correlation $r_{\rm rb}$ (positive = PSO-VNS better).}
\label{tab:wtl}
\scriptsize\setlength{\tabcolsep}{2.5pt}
\begin{tabular}{llcccccccc}
\toprule
DS & $r$ (m) & Cases & PSO & SSA-VNS & SSA & LX-SSA & DE & VNS & MS-SLSQP \\
\midrule
I & 500 & 9 & 0/8/1 & 6/3/0 & 8/1/0 & 8/1/0 & 7/2/0 & 6/3/0 & 7/2/0 \\
I & 750 & 11 & 1/8/2 & 9/2/0 & 10/1/0 & 10/1/0 & 9/2/0 & 8/3/0 & 10/1/0 \\
I & 1000 & 14 & 3/11/0 & 9/5/0 & 12/2/0 & 12/2/0 & 12/2/0 & 12/2/0 & 12/2/0 \\
II & 500 & 9 & 3/5/1 & 7/2/0 & 7/2/0 & 8/1/0 & 7/2/0 & 7/1/1 & 8/1/0 \\
II & 750 & 11 & 2/9/0 & 8/3/0 & 10/1/0 & 10/1/0 & 9/2/0 & 8/3/0 & 9/2/0 \\
II & 1000 & 14 & 5/9/0 & 10/4/0 & 12/2/0 & 12/2/0 & 12/2/0 & 10/4/0 & 13/1/0 \\
\midrule
All & & 68 & 14/50/4 & 49/19/0 & 59/9/0 & 60/8/0 & 56/12/0 & 51/16/1 & 59/9/0 \\
$\tilde r_{\rm rb}$ & & & +0.08 & +0.72 & +0.98 & +0.97 & +1.00 & +0.84 & +0.98 \\
\bottomrule
\end{tabular}
\end{table}


Over the eight methods (Table~\ref{tab:friedman68}), the Friedman test \NFriedVerb{} the hypothesis of equal performance ($\chi^2_F=\NFriedChi$, $p=\NFriedP$; Iman--Davenport $F_F=\NImanF$, $p=\NImanP$). PSO-VNS has the best average rank (\NRankPSOVNS), alone ranks first in \NSoleBestPSOVNS{} cases, and stays first when the final layouts are re-evaluated with 5$^\circ$ or 1$^\circ$ direction sub-bins (1$^\circ$: \NFRankPSOVNSOne; Section~\ref{sec:robustness}). In the Holm-adjusted post hoc tests ($p_z$ in Table~\ref{tab:friedman68}; tied ranks make them conservative), PSO-VNS ranks significantly better than \NPostHocSigList{}, but not than \NPostHocNonSigList{} ($p_{\rm Holm}=\NPostHocNonSigP$). The rank order and these conclusions are unchanged for thresholds of \NXThrList{} feasible runs; significant comparisons with PSO-VNS favor the same method in each of the \NXClusters{} clusters, and leaving out one cluster changes none of them (Tables~\ref{S-tab:X-threshold}--\ref{S-tab:X-loco}).
% CHECK-FINAL [C01]: main.friedman.best_ranked == PSOBV; [C02]: p < 0.05 (verb generated: \NFriedVerb); [C03]: p_holm_vs_focus: only PSOC >= 0.05
% CHECK-DIR [F06]: PSO-VNS best average rank with 5-deg and 1-deg sub-bins (\NFRankPSOVNSOne)
% CHECK-EXTRA [X02][X03][X05]-[X09] six thresholds 1-30; [X11] clusters 6/6 same direction for comparisons with PSO-VNS; [X14] LOCO changes only \NXLocoChanged (not a PSO-VNS comparison), PSO-VNS always best

Against PSO, PSO-VNS is significantly better in \NWtlPSOW{} cases, tied in \NWtlPSOT{} and worse in \NWtlPSOL{} (run level, Table~\ref{tab:wtl}). The case means do not differ significantly ($\overline{\Delta L}=\NCmPSOVNSvsPSODL$~pp, $p=\NCmPSOVNSvsPSOP$; Hodges--Lehmann estimate \NXHLPSOVNSvsPSO~pp). Whether the difference is negligible at the post hoc margin $\pm\NEqMargin$~pp depends on the level of inference (Table~\ref{S-tab:X-equiv-levels}): on this fixed benchmark (seed level), PSO-VNS has a small significant advantage inside the margin (90\% CI \NXEqSeedPSOVNSvsPSOCI); over exchangeable cases the two are equivalent (\NXEqCasePSOVNSvsPSOCI, TOST $p=\NXEqCasePSOVNSvsPSOP$, minimal margin \NXEqCasePSOVNSvsPSOMin); over farm clusters, equivalence is not shown (CR2 \NXEqClustPSOVNSvsPSOCI, minimal margin \NXEqClustPSOVNSvsPSOMin). In the Bayesian signed-rank test, practical equivalence is the most probable outcome (posterior probability \NBayPSOVNSvsPSORope; posterior means \NXBayMeanPSOVNSvsPSOLeft/\NXBayMeanPSOVNSvsPSORope/\NXBayMeanPSOVNSvsPSORight{} for PSO-VNS better by more than the margin/equivalent/PSO better). The margin corresponds to \NXMarginMWhTurbDsI{} (Data Set~I) and \NXMarginMWhTurbDsII~MWh/yr (Data Set~II) per turbine, the mean gain of PSO-VNS to \NXEnPSOVNSvsPSOMWhTurb~MWh/yr per turbine (\NXEnPSOVNSvsPSOPct\% of AEP). Spread and worst run do not differ ($p=\NSdP$, \NWorstP).
% CHECK-FINAL [C04]: PSOC: W >= L, mean loss PSOBV < PSOC, case-mean p >= 0.05; [C49]: case-level 90% CI inside +-m, P(rope most probable) printed via \NBayPSOVNSvsPSORope; [C52]: SD p and worst-run p >= 0.05
% CHECK-EXTRA [X39] seed and case level equivalent, cluster level not; [X40] seed-level CI excludes 0 inside m; [X41] minimal margins seed < case < m < cluster; [X48] Bayesian wording + posterior means; [X50] HL near 0; [X56] margin 4.1/2.1 MWh/yr per turbine, mean gain 0.2; [X25] energy < 0.05 % AEP

The average hides case-level differences in both directions: \NXHetBeyondPSOVNS{} cases favor PSO-VNS by more than the margin (by up to \NXHetMaxPSOVNS~pp) and \NXHetBeyondPSO{} favor PSO (up to \NXHetMaxPSO~pp; Table~\ref{S-tab:X-heterogeneity}). The two are equivalent for $N<10$ (\NXEqSmallN{} cases, 90\% CI \NXEqSmallCI) but not for $N\ge10$ (\NXEqLargeMean~pp, \NXEqLargeCI) or on the \NXEqNontrivialN{} cases with a wake loss of at least \NXTrivialLoss~pp (\NXEqNontrivialCI). They follow a data set $\times$ density crossover (density $\phi=\NXDensDef$; stated after review, exploratory; interaction $p=\NXDensInterP$): as $\phi$ grows, PSO-VNS gains in Data Set~II (slope \NXDensSlopeDsII~pp, $p=\NXDensPDsII$) and PSO in Data Set~I (\NXDensSlopeDsI~pp, $p=\NXDensPDsI$). In the post hoc subgroup Data Set~II, $N\ge10$, PSO-VNS has the lower mean wake loss in \NCmPSOVNSvsPSOdsIILargeWins{} of \NCmPSOVNSvsPSOdsIILargeN{} cases ($p_{\rm Holm}\le\NXHolmPSOVNSvsPSOdsIILargeP$ over all \NXMultN{} case-mean Wilcoxon tests; \NXEnPSOVNSvsPSOdsIILargePct\% of AEP); all \NXHetDsINPSO{} Data Set~I cases beyond the margin favor PSO. With 1$^\circ$ sub-bins, PSO-VNS is slightly but significantly better than PSO ($\overline{\Delta L}=\NFPairMeanOne$~pp, 90\% CI \NFPairCIOne; Table~\ref{S-tab:F-equiv}).
% CHECK-EXTRA [X43] 11/8 beyond +-m, max 0.41/0.27; [X44] N < 10 equivalent, N >= 10 not; [X45] non-trivial cases not equivalent; [X47] data set x density crossover (exploratory, stated after review); [X19] DS II N >= 10 survives all-family Holm ("<=", R5 O13); [X27] DS II N >= 10 ~0.2 % AEP
% CHECK-FINAL [C32]: DS II N >= 10 case-mean wins > losses
% CHECK-DIR [F10]: 1-deg PSO-VNS - PSO CI below 0, not equivalent; [F14]: DS II N >= 10 advantage survives at 1 deg

Against SSA-VNS, PSO-VNS is significantly better in \NWtlSSAVNSW{} cases and never worse ($p_{\rm Holm}=\NPWSSAVNS$). Against SSA, LX-SSA, DE, VNS and MS-SLSQP (untuned DE; finite-difference gradients), it is significantly better in most cases; \NLossesPhrase. PSO-VNS is more reliable in reaching feasibility on the Horns Rev~1 block with random starts (\NHRFeasPSOVNS{} vs.\ \NHRFeasPSO{}, McNemar $p=\NXHRMcNemarP$; Section~\ref{sec:hornsrev-site}) and ranks best on the six largest cases at all three budgets (exploratory; Section~\ref{sec:budget}). The leading methods depend on the conditions: from feasible starts, \NFbFeasBestRank{} ranks first (Section~\ref{sec:feasinit}); under the Gaussian wake, \NBestGauss{} ranks narrowly ahead of PSO-VNS with 15$^\circ$ bins but not with 1$^\circ$ sub-bins (Section~\ref{sec:gaussian}).
% CHECK-FINAL [C06]: wtl.SSABV.L == 0; [C07]: better than SSA, LX-SSA, DE, VNS, MS-SLSQP in most cases; [C08]/[C20]: Horns Rev PSO-VNS 30/30, PSO < 30/30; [C29]: six largest cases best at all budgets; [C28]: feasible-start BVNS best; [C22]: Gauss 15 deg PSOC first, PSOBV second
% CHECK-EXTRA [X52] McNemar (one site, random starts); CHECK-DIR [F06] best rank 15/5/1 deg; [F09] Gaussian: PSO first at 15 deg, PSO-VNS at 1 deg

\begin{figure*}[!t]
\centering
\pipefig{0.7\textwidth}{figures_mpce/ablation_convergence.pdf}{6.6cm}
\caption{Convergence of the \NAblNVar{} component-analysis variants (Table~\ref{tab:ablation}) for the largest turbine count of each farm: median best feasible wake loss over 30 runs versus evaluations. PSO-VNS and PSO coincide up to 3,030 evaluations (dotted line), where PSO-VNS switches to VNS. Curves of the main comparison: Fig.~\ref{S-fig:S-conv-max} of the supplementary material.}
\label{fig:conv-max}
\label{fig:ablation}
\end{figure*}
% CHECK (pipeline count): \NAblNVar = variants of tab:ablation = curves of ablation_convergence.pdf; [C58] ranks

The wake loss grows with $N$, falls as the farm becomes larger, and is higher for the broader wind rose of Data Set~II (Section~\ref{S-sec:S-figs}). PSO-VNS ends feasible in all runs (DE: \NFeasDE\%). In the densest cases (500~m, $N=10$), the global best of the PSO phase is feasible at the switch in only \NDenseSwitchFeas\% of the runs, but all runs are feasible after the VNS phase.
% CHECK-FINAL [C42]: case-mean wake loss grows with N, falls with r, DS II >= DS I; [C08]: feasible_pct.PSOBV == 100; [C24]: dense_500_10 feasible at switch < 100

After the switch, both keep improving (Fig.~\ref{fig:conv-max}): the VNS phase removes on average \NPhTwoMean\% (median \NPhTwoMedian\%) of the wake loss left after the PSO phase and repairs \NPhTwoMadeFeas{} runs infeasible at the switch, whereas continuing PSO for the same 3,000 evaluations removes \NPsoContMean\% (median \NPsoContMedian\%).
% CHECK-FINAL [C10]: phase2 PSOBV.mean >= PSOC_continued.mean > 0 ("both keep improving")
